A lo largo de esta práctica, se logró observar de manera experimental la importancia 
crítica del preprocesamiento y limpieza de datos, imágenes y señales dentro de las 
tareas de Reconocimiento de Patrones. 

En la primera sección, el análisis exploratorio y de correlación del conjunto de datos 
Iris nos permitió identificar dependencias positivas fuertes entre las medidas morfológicas 
(como la longitud y ancho de los pétalos), lo cual es fundamental para el posterior 
entrenamiento de modelos de clasificación geométrica.

En la segunda sección, comprobamos que la aplicación de filtros de suavizado (destacando 
el filtro Gaussiano como el resultado óptimo) es un paso indispensable antes de aplicar 
técnicas de extracción de características como el operador de Sobel/Prewitt o la 
Transformada de Hough. El ruido en la imagen original generaba falsos positivos masivos 
en la detección de bordes y círculos, los cuales fueron mitigados eficazmente mediante 
el filtrado.

Finalmente, en el procesamiento de señales de ECG, pudimos comprobar la efectividad
matemática de las transformaciones de datos. La señal ruidosa ocultaba casi por 
completo las ondas de baja amplitud (P y T) del latido. Sin embargo, demostramos que 
tanto el truncamiento espectral en el dominio de la frecuencia (vía FFT/IFFT) como el 
uso de un filtro IIR en el dominio del tiempo (Butterworth) son altamente efectivos 
para atenuar la interferencia de alta frecuencia y recuperar la integridad de la 
morfología original de la señal cardiovascular.

\QED